\xmpheader 4/{Spacing, rules, and boxes}
Here's an example of a ``description list''. \par
% Call the indentation for descriptions \descindent, to avoid repetitive
% constructs, we'' define a macro
\newdimen\descindent \descindent = 8pc
\def\describe#1{\par\noindent \llap{\hbox to \descindent{\bf #1\hfil}}}
{\leftskip = \descindent \parskip = 0.5 \baselineskip
% Move the description to the left of the paragraph.
\describe{Queen of Hearts}An ill-tempered woman, prone to saying ``Off with his head!'' 
at the slightest provocation.
\describe{Cheshire Cat}A cat with an enormous smile that Alice found in a tree.
\describe{Mock Turtles}A lachrymose creature, quite a storyteller, who was a companion 
to the Cryphon.  Reputedly the principal ingredient of Mock Turtle Soup.
\par}
\bigskip\hrule\bigskip % A line with vertical space around it.
Here's an example of some words in a ruled box, just as Lewis Carroll wrote them:
\medskip
% Put 8pt of space between the text and the surrounding rules.
\hbox{\vrule\vbox{\hrule
   \hbox spread 8 pt{\hfil\vbox spread 8pt {\vfil
       \hbox{Who would not give all else for twop}%
       \hbox{ennyworth only of Beautiful Soup?}%
   \vfil}\hfil}
\hrule}\vrule}%
\medskip
\smallskip
\line{\hfil\hbox to 3 in{\leaders\hbox{ * }\hfil}\hfil}
\bigskip
\line{\hskip -4pt\vrule\hfil\vbox{Here we've gotten the effect of a revision bar on the material 
in this paragrah.  The revision bar might indicate a change.}}

You use the Unix command {\tt tex $file$} to evoke the \TeX\ program; which translate your \TeX\ 
source into {\tt $file$.dvi}. Then {\tt dvips $file$} produces {\tt $file$.ps}, which you can
send to printer. You can view the {\tt dvi} file by {\tt xdvi} and {\tt ps} file by {\tt ghostview}.
